\documentclass[a4paper,11pt]{article}
\usepackage[a4paper,margin=2cm]{geometry}
\usepackage[T1]{fontenc}
\usepackage[utf8]{inputenc}
\usepackage[ngerman,english]{babel}
\usepackage{lmodern,microtype,xcolor,parskip,enumitem}
\usepackage[most]{tcolorbox}
\usepackage{titlesec}
\usepackage[hidelinks]{hyperref}
\definecolor{HeaderBlueDark}{RGB}{70,110,170}
\definecolor{RowBlueLight}{RGB}{235,243,255}
\definecolor{RowBlueDark}{RGB}{220,233,250}
\definecolor{SoftGray}{RGB}{90,90,90}
\setlength{\parindent}{0pt}
\setlist[itemize]{leftmargin=1.4em,itemsep=0.35em,topsep=0.35em}
\titleformat{\section}[block]{\Large\bfseries\color{HeaderBlueDark}}{}{0pt}{}
\newtcolorbox{exbox}{enhanced,breakable,colback=RowBlueLight,colframe=RowBlueDark,boxrule=0.4pt,arc=1.5mm,left=1.2mm,right=1.2mm,top=1mm,bottom=1mm,title={Beispiele \textnormal{(Examples)}},fonttitle=\bfseries,colbacktitle=RowBlueDark,coltitle=black}
\NewDocumentEnvironment{grammarentry}{mm+b}{\begin{tcolorbox}[enhanced,breakable,colback=white,colframe=HeaderBlueDark,title={#1\hfill\normalfont\footnotesize #2},colbacktitle=HeaderBlueDark,coltitle=white,fonttitle=\bfseries]#3\end{tcolorbox}}{}
\newcommand{\GermanText}[1]{\textbf{Deutsch: }#1\par}
\newcommand{\EnglishText}[1]{{\small\color{SoftGray}\textbf{English: }#1\par}}
\newcommand{\ExampleItem}[2]{\item \textbf{#1}\\{\small\color{SoftGray}#2}}
\title{Die Modalpartikel \glqq halt\grqq\\{\large\textnormal{The modal particle halt}}}
\date{}
\begin{document}
\maketitle
\tableofcontents
\section{Bedeutung und Wirkung \textnormal{(Meaning and effect)}}
\begin{grammarentry}{halt}{Modalpartikel / Abtönungspartikel}
\GermanText{Halt zeigt an, dass man eine Tatsache als gegeben oder schwer veränderbar akzeptiert.}
\EnglishText{Halt presents a situation as given or difficult to change.}
\end{grammarentry}
\section{Grammatische Rolle und Satzarten \textnormal{(Grammar and sentence moods)}}
\begin{grammarentry}{Verwendung}{Satzart und Kontext}
\GermanText{Besonders häufig in mündlichen Aussagen und beiläufigen Ratschlägen, oft mit Resignation.}
\EnglishText{Especially common in spoken statements and casual advice, often with a resigned tone.}
\end{grammarentry}
\begin{grammarentry}{Form und Stellung}{Syntax}
\GermanText{Als Modalpartikel ist halt unflektierbar, fordert keinen Kasus und steht normalerweise im Mittelfeld. Allein im Vorfeld eines normalen Aussagesatzes kann eine Modalpartikel nicht stehen. Die Betonung und der Kontext beeinflussen die Wirkung.}
\EnglishText{As a modal particle, halt is invariable, does not govern case and normally appears in the middle field. It cannot occupy the prefield of a normal statement on its own. Stress and context affect its pragmatic force.}
\end{grammarentry}
\section{Typische Beispiele \textnormal{(Typical examples)}}
\begin{exbox}\begin{itemize}
\ExampleItem{Das ist halt so.}{That's just how it is.}
\ExampleItem{Dann geh halt allein.}{Then just go by yourself.}
\ExampleItem{Er hat halt keine Zeit.}{He simply doesn't have time.}
\ExampleItem{Wir müssen halt warten.}{We'll just have to wait.}
\end{itemize}\end{exbox}
\section{Abgrenzung zu anderen Verwendungen \textnormal{(Contrast with other uses)}}
\begin{grammarentry}{Wortarten unterscheiden}{Nicht immer Modalpartikel}
\GermanText{Halt als Nomen (der Halt) oder als Imperativ von halten hat eine andere Wortart.}
\EnglishText{Halt as a noun (support/stop) or as an imperative of halten belongs to another word class.}
\end{grammarentry}
\begin{exbox}\begin{itemize}
\ExampleItem{Der Bus macht hier Halt.}{The bus stops here.}
\ExampleItem{So ist es eben.}{That's just the way it is.}
\end{itemize}\end{exbox}
\section{Lernhinweis \textnormal{(Learning tip)}}
\GermanText{Eine Modalpartikel verändert vor allem die Gesprächsabsicht. Versuchen Sie nicht, jedem Vorkommen eine feste englische Eins-zu-eins-Übersetzung zu geben. Achten Sie auf Satzart, Tonfall und Situation.}
\EnglishText{A modal particle mainly affects the speaker's intention. Do not try to give it one fixed English translation in every context. Pay attention to sentence mood, intonation and situation.}
\section{Quellen \textnormal{(Sources)}}
\begin{itemize}
\item IDS Mannheim: \url{https://grammis.ids-mannheim.de/terminologie/2}
\item Duden: \url{https://www.duden.de/sprachwissen/sprachratgeber/Modalpartikeln}
\end{itemize}
\end{document}
